\documentclass{article}
\input{../../common_latex/import_include}
\input{../../common_latex/import_authors}
%-----------------------------------------------
% title
% 
\title{
	\Huge \textbf{NR Access} \\[1em]
}
\author{}
\date{}

%-----------------------------------------------
% document
%
\begin{document}

\maketitle

\tableofcontents
\newpage

\section{Terminology}
\input{../../common_latex/import_terms}

\section{介绍}
本节课程主要介绍NR接入流程。



\section{PRACH接入}
\subsection{工程经验}
UE侧的处理还会延迟约5个TA，物理层针对Preamble检测预留了2个TA的保护。如果考虑工程经验，覆盖范围还需要缩减。

\subsection{最大小区半径}
首先我们需要根据小区半径来选择PRACH Format来发送。
\begin{itemize}
\item Format 0$\to$3（仅支持FR1，详见TS38211 Table 6.3.3.2-2, Table 6.3.3.2-3, Table 6.3.3.2-4）\\
TS38211 5.3.2 OFDM baseband signal generation for PRACH原文如下"$N_{CP,l}^{RA}=N_{CP}^{RA}+n∙16k$ where for $\Delta f_{RA}={1.25,5}kHz$, n=0"\cite{38211}。这个说明对于format0$\to$3，原始的CP长度就是PRACH的CP长度。对于format0$\to$3，调度的slot是按照15kHz调度的\footnote{TS38211 6.3.3.2 Mapping to physical resources规定FR1按照15kHz计算slot数量，而format0$\to$3只占用FR1，且NR的调度单元是slot，所以一旦调度PRACH format 0$\to$3就会按照15kHz slot连续调度}。

为了避免PRACH自我干扰和干扰到其他调度单元，UE和BS空口传输时间$T_{prop}$必须满足
\begin{equation}
    min\{N_{CP}^{RA}, N_{seq}, N_{gap}\} \geq 2T_{prop}\footnotemark
\end{equation}
\footnotetext{这里是2$T_{prop}$：首先UE的时间窗口比BS的时间窗口晚了$T_{prop}$；为了克服UE发送中的时延，$T_{prop}$，UE相对自己的时间窗需要提前$T_{prop}$；需要相对BS的时间窗口提前2$T_{prop}$\label{fn:PRACH-2Tprop_reason}}
其中$N_{seq}$是PRACH中一个preamble的长度。

\begin{itemize}
\item Format 0: $N_{CP}^{RA} + N_u$=3168k+24576k=27744$k$=0.096875(ms), $N_{gap}$=1-0.096875=2976k\\
\[
r_{cell} = N_{gap}\cdot 3e8/2=14.5312(km)
\]
\item Format 1: $N_{CP}^{RA} + N_u$=21024k+2$\cdot$24576k=70176$k$=2.2844(ms), $N_{gap}$=3-2.2844=21984k
\[
r_{cell} = N_{CP}^{RA}\cdot 3e8/2=102.6562(km)
\]
\item Format 2: $N_{CP}^{RA} + N_u$=4688k+4$\cdot$24576k=102992$k$=3.3526(ms), $N_{gap}$=3.5-3.3526=4528k
\[
r_{cell} = N_{gap}\cdot 3e8/2=22.1094(km)
\]
\item Format 3: $N_{CP}^{RA} + N_u$=3168k+4$\cdot$24576k=27744$k$=0.096875(ms), $N_{gap}$=1-0.096875=2976k
\[
r_{cell} = N_{gap}\cdot 3e8/2=14.5312(km)
\]
\end{itemize}

\item Format A1$\to$A3支持\{15, 30, 60, 120, 480, 960\}kHz（没有$N_{gap}$）
\begin{itemize}
\item Format A1: $N_{CP}^{RA} + N_u$=$(288k+2\cdot 2048k)\cdot 2^{-\mu}$=$2\cdot(144k+2048k)\cdot 2^{-\mu}$, $N_{gap}$=0
\[
r_{cell} = N_{CP}^{RA}\cdot 3e8/2=\begin{cases}
1.4062(km) & scs=15kHz \\
0.7031(km) & scs=30kHz \\
0.3516(km) & scs=60kHz \\
0.1758(km) & scs=120kHz \\
0.0439(km) & scs=480kHz \\
0.0220(km) & scs=960kHz
\end{cases}
\]
\item Format A2: $N_{CP}^{RA} + N_u$=$(576k+4\cdot 2048k)\cdot 2^{-\mu}$=$4\cdot(144k+2048k)\cdot 2^{-\mu}$, $N_{gap}$=0
\[
r_{cell} = N_{CP}^{RA}\cdot 3e8/2=\begin{cases}
2.8125(km) & scs=15kHz \\
1.4062(km) & scs=30kHz \\
0.7031 (km) & scs=60kHz \\
0.3516(km) & scs=120kHz \\
0.0879(km) & scs=480kHz \\
0.0439(km) & scs=960kHz
\end{cases}
\]
\item Format A3: $N_{CP}^{RA} + N_u$=$(864k+6\cdot 2048k)\cdot 2^{-\mu}$=$6\cdot(144k+2048k)\cdot 2^{-\mu}$, $N_{gap}$=0
\[
r_{cell} = N_{CP}^{RA}\cdot 3e8/2=\begin{cases}
4.2188(km) & scs=15kHz \\
2.1094(km) & scs=30kHz \\
1.0547(km) & scs=60kHz \\
0.5273(km) & scs=120kHz \\
0.1318(km) & scs=480kHz \\
0.0659(km) & scs=960kHz
\end{cases}
\]
\end{itemize}

\item Format B1$\to$B4, Format C0,C2支持\{15, 30, 60, 120, 480, 960\}kHz（由于每次重复都占用一个OFDM symbol，那么所有slot中CP时间扣除$N_{CP}^{RA}$后就是$N_{gap}$）
\begin{itemize}
\item Format B1, $N_{CP}^{RA}$=$216k\cdot 2^{-\mu}$, $N_{gap}$=$72k\cdot 2^{-\mu}$
\[
r_{cell} = N_{gap}\cdot 3e8/2=\begin{cases}
0.3516(km) & scs=15kHz \\
0.1758(km) & scs=30kHz \\
0.0879(km) & scs=60kHz \\
0.0439(km) & scs=120kHz \\
0.0110(km) & scs=480kHz \\
0.0055(km) & scs=960kHz
\end{cases}
\]
\item Format B2, $N_{CP}^{RA}$=$360k\cdot 2^{-\mu}$, $N_{gap}$=$216k\cdot 2^{-\mu}$
\[
r_{cell} = N_{gap}\cdot 3e8/2=\begin{cases}
1.0547(km) & scs=15kHz \\
0.5273(km) & scs=30kHz \\
0.2637(km) & scs=60kHz \\
0.1318(km) & scs=120kHz \\
0.0330(km) & scs=480kHz \\
0.0165(km) & scs=960kHz
\end{cases}
\]
\item Format B3, $N_{CP}^{RA}$=$504k\cdot 2^{-\mu}$, $N_{gap}$=$360k\cdot 2^{-\mu}$
\[
r_{cell} = N_{gap}\cdot 3e8/2=\begin{cases}
1.7578(km) & scs=15kHz \\
0.8789(km) & scs=30kHz \\
0.4395(km) & scs=60kHz \\
0.2197(km) & scs=120kHz \\
0.0549(km) & scs=480kHz \\
0.0275(km) & scs=960kHz
\end{cases}
\]
\item Format B4, $N_{CP}^{RA}$=$936k\cdot 2^{-\mu}$, $N_{gap}$=$792k\cdot 2^{-\mu}$
\[
r_{cell} = N_{gap}\cdot 3e8/2=\begin{cases}
3.8672(km) & scs=15kHz \\
1.9336(km) & scs=30kHz \\
0.9668(km) & scs=60kHz \\
0.4834(km) & scs=120kHz \\
0.1208(km) & scs=480kHz \\
0.0604(km) & scs=960kHz
\end{cases}
\]
\item Format C0, $N_{CP}^{RA}$=$1240k\cdot 2^{-\mu}$, $N_{gap}$=$1096k\cdot 2^{-\mu}$
\[
r_{cell} = N_{CP}^{RA}\cdot 3e8/2=\begin{cases}
5.3516(km) & scs=15kHz \\
2.6758(km) & scs=30kHz \\
1.3379(km) & scs=60kHz \\
0.6689(km) & scs=120kHz \\
0.1672(km) & scs=480kHz \\
0.0836(km) & scs=960kHz
\end{cases}
\]
\item Format C2, $N_{CP}^{RA}$=$2048k\cdot 2^{-\mu}$, $N_{gap}$=$2912k\cdot 2^{-\mu}$
\[
r_{cell} = N_{CP}^{RA}\cdot 3e8/2=\begin{cases}
10(km) & scs=15kHz \\
5(km) & scs=30kHz \\
2.5(km) & scs=60kHz \\
1.25(km) & scs=120kHz \\
0.3125(km) & scs=480kHz \\
0.1562(km) & scs=960kHz
\end{cases}
\]
\end{itemize}
\end{itemize}

\subsubsection{计算当前小区适配的$N_{cs}$}
上面给出的最大小区半径是$N_{cs}=0$，如果小区半径在最大半径以内，则可以计算$N_{cs}$最小值以满足最大半径覆盖范围(如果考虑定时误差、多径时延冗余、安全冗余还需要增加$T_{sec}$)。

\begin{figure}[htbp]
    \centering
    \includegraphics[width=0.5\textwidth, keepaspectratio]{lect001_access_tdd_guard}
    \caption{TDD Guad时间}
    \label{fig:example}
\end{figure}
\[
N_{cs} = L_{RA}\lfloor \frac{2r_{cell}/c+T_{sec}}{T_{seq}} \rfloor
\]


\subsection{小区半径对TDD Slot结构的限制}
小区半径此时已经确认，TDD传输时DL$\to$UL切换时，需要预留保护间隔（GAP），确保上下行隔离，不存在交叠。如下图所示，GAP的长度要求大于下行延时量和上行提前量之和，即$2T_{prop}+N_{TA,offset}$，其中$T_{prop}$表示小区边缘用户的信号接收时延。
\[
T_{guard} \geq N_{TA,offset} + 2T_{prop}\textsuperscript{\ref{fn:PRACH-2Tprop_reason}}
\]
其中$T_{prop}$就是小区半径导致的传输延时。
\begin{longtable}{|p{0.15\textwidth}|p{0.15\textwidth}|p{0.15\textwidth}|}
\caption{PRACH Long Format $T_{guard} $}\\
\hline
 & \multicolumn{1}{|c|}{$N_{TA,offset}=25600$} & \multicolumn{1}{|c|}{$N_{TA,offset}=39936$} \\ \hline
Format 0 & 3376k & 3600k \\ \hline
Format 1 & 22384k & 22608k \\ \hline
Format 2 & 4928k & 5152k \\ \hline
Format 3 & 3376k & 3600k \\ \hline
\end{longtable}

\begin{longtable}{|p{0.15\textwidth}|p{0.15\textwidth}|p{0.15\textwidth}|p{0.15\textwidth}|}
\caption{PRACH Short Format $T_{guard} $}\\
\hline
& \multicolumn{1}{|c|}{$N_{TA,offset}=25600$} & \multicolumn{1}{|c|}{$N_{TA,offset}=39936$} & \multicolumn{1}{|c|}{$N_{TA,offset}=13792$} \\ \hline
Format A1 & 288k$\cdot 2^{-\mu}$+400 & 288k$\cdot 2^{-\mu}$+624 & 288k$\cdot 2^{-\mu}$+215.5 \\ \hline
Format A2 & 576k$\cdot 2^{-\mu}$+400 & 576k$\cdot 2^{-\mu}$+624 & 576k$\cdot 2^{-\mu}$+215.5 \\ \hline
Format A3 & 864k$\cdot 2^{-\mu}$+400 & 864k$\cdot 2^{-\mu}$+624 & 864k$\cdot 2^{-\mu}$+215.5 \\ \hline
Format B1 & 72k$\cdot 2^{-\mu}$+400 & 72k$\cdot 2^{-\mu}$+624 & 72k$\cdot 2^{-\mu}$+215.5 \\ \hline
Format B2 & 216k$\cdot 2^{-\mu}$+400 & 216k$\cdot 2^{-\mu}$+624 & 216k$\cdot 2^{-\mu}$+215.5 \\ \hline
Format B3 & 360k$\cdot 2^{-\mu}$+400 & 360k$\cdot 2^{-\mu}$+624 & 360k$\cdot 2^{-\mu}$+215.5 \\ \hline
Format B4 & 792k$\cdot 2^{-\mu}$+400 & 792k$\cdot 2^{-\mu}$+624 & 792k$\cdot 2^{-\mu}$+215.5 \\ \hline
Format C0 & 1096k$\cdot 2^{-\mu}$+400 & 1096k$\cdot 2^{-\mu}$+624 & 1096k$\cdot 2^{-\mu}$+215.5 \\ \hline
Format C2 & 2048k$\cdot 2^{-\mu}$+400 & 2048k$\cdot 2^{-\mu}$+624 & 2048k$\cdot 2^{-\mu}$+215.5 \\ \hline
\end{longtable}



\paragraph{$N_{TA,offset}$}\mbox{}\\
\begin{itemize}
\item $N_{TA,offset}$由高层RRC n-TimingAdvanceOffset（出现在ServingCellConfigCommon和\\
ServingCellConfigCommonSIB 中），有$n0, n25600, n39936$三个值（表示对应的数量的tc）\cite{38331}。
\item 高层RRC n-TimingAdvanceOffset没有配置时，UE采用38133给的默认值（TS38213 4.2Transmission timing adjustments的原文是"If the UE is not provided n-TimingAdvanceOffset for a serving cell, the UE determines a default value NTA,offset of the timing advance offset for the serving cell as described in [10, TS 38.133]"）\cite{38213}。
\begin{itemize}
\item TS38133 7.1.1	Introduction原文"$N_{TA}$ offset for first TAG and second TAG are as defined in n-TimingAdvanceOffset and n-TimingAdvanceOffset2 respectively, if configured, otherwise it is $N_{TA,offset}$ defined in table 7.1.2-2."\cite{38133}。这说明缺省时，采用Table 7.1.2-2配置。

\begin{longtable}{|p{0.8\textwidth}|p{0.19\textwidth}|}
\caption{TS38133 Table 7.1.2-2\cite{38133}}\\
\hline
\multicolumn{1}{|c|}{Frequency range and band of cell used for uplink transmission} & \multicolumn{1}{c|}{$N_{TA,offset}$(Unit: $T_c$)} \\ \hline
FR1 FDD or TDD band with neither E-UTRA–NR nor NB-IoT–NR coexistence case & 25600 \\ \hline
FR1 FDD band with E-UTRA–NR and/or NB-IoT–NR coexistence case & 0 \\ \hline
FR1 TDD band with E-UTRA–NR and/or NB-IoT–NR coexistence case & 39936 \\ \hline
FR2 & 13792 \\ \hline
\end{longtable}

\end{itemize}

\end{itemize}




\section{故障排查}


\subsection{Background}



\section{Time Frequency Domain}
todo

\section{PRACH}
\subsection{Root Sequence}
\begin{itemize}
\item 对于非限制集，所有的根序列u_i都是可用的（见TS38211 表6.3.3.1-2, 6.3.3.1-3, 6.3.3.1-4A, 6.3.3.1-4B ）

\item 对于限制集A和B，我们需要检查当前根序列u_i中是否有可用的preamble\\
计算最小非负整数$q$, $(qu)mod L_{RA}=1$
\[
d_u = \begin{cases}
q & 0 \leq q \leq L_{RA}/2\\
L_{RA} - q & q > L_{RA}/2\\
\end{cases}
\]
\begin{itemize}
\item 限制集A需要满足$N_{cs} \leq d_u \leq (L_{RA}−N_{cs})/2$
\item 限制集B需要满足3个条件之一
\[
\begin{split}
N_{cs} &\leq d_u \leq (L_{RA} - N_{cs})/4 \\
(L_{RA}+N_{cs})/4 &\leq d_u \leq (L_{RA} - N_{cs})/3 \\
(L_{RA}+N_{cs})/3 &\leq d_u \leq (L_{RA} - N_{cs})/2
\end{split}
\]
\end{itemize}
\end{itemize}
\subsubsection{$q$的由来}
从根序列索引 $u$ 推导到其模逆元 $q$本质上是利用了高斯和（Gauss Sum）的性质以及离散傅里叶变换（DFT）的二次相位保持特性。
\begin{align}
x_u(i) &= e^{-j\frac{\pi ui(i+1)}{L_{RA}}} \\
X_u(k) &= \sum_{i=0}^{L_{RA}-1} e^{-j\frac{\pi ui(i+1)}{L_{RA}}}e^{-j\frac{2\pi ki}{L_{RA}}} \\
&= \sum_{i=0}^{L_{RA}-1} e^{-j\frac{\pi}{L_{RA}} (ui^2 + ui + 2ki)   }\label{prach-Xu_orig}
\end{align}
其中$e^{-j\frac{2\pi ki}{L_{RA}}}$可以进行重新推导，设$qu=mL_{RA}+1$
\begin{align}
e^{-j\frac{2\pi ki}{L_{RA}}} &= e^{-j\frac{2\pi ki + 2\pi ki mL_{RA}}{L_{RA}}} \\
&= e^{-j\frac{2\pi ki(mL_{RA} +1)}{L_{RA}}} \\
&= e^{-j\frac{2\pi ki qu}{L_{RA}}}\label{prach-1_qu} 
\end{align}
把\eqref{prach-1_qu}带入\eqref{prach-1_qu}得到
\begin{align}
X_u(k) &= \sum_{i=0}^{L_{RA}-1} e^{-j\frac{\pi}{L_{RA}} (ui^2 + ui + 2kiuq)   } \\
&= \sum_{i=0}^{L_{RA}-1} e^{-j\frac{\pi}{L_{RA}} u(i^2 + i + 2kqi) } \\
&= \sum_{i=0}^{L_{RA}-1} e^{-j\frac{\pi}{L_{RA}} u(i^2 + 2\frac{1 + 2kq}{2}i + \frac{(1 + 2kq)^2}{4} - \frac{(i + 2kq)^2}{4})} \\
&= \sum_{i=0}^{L_{RA}-1} e^{-j\frac{\pi}{L_{RA}}u(i + \frac{1 + 2kq}{2})^2} e^{j\frac{\pi}{L_{RA}}u\frac{(1 + 2kq)^2}{4}} \\
&= \underbrace{e^{j\frac{u\pi(1 + 2kq)^2}{4L_{RA}}}}_{phase}\underbrace{\sum_{i=0}^{L_{RA}-1} e^{-j\frac{\pi}{L_{RA}}u(i + \frac{1 + 2kq}{2})^2}}_{Quad-Gauss-Sum}
\end{align}
Let's look at the Quad-Gauss-Sum part
\begin{align}
dg
\end{align}



\subsection{}
%5G 的时间单位$t_c=\frac{1}{\Delta f_{\max}\cdot 4096}$, 5G 的FFT size最大是4096（）
\begin{equation}
s_l^{p,\mu} = \sum_{k=0}^{L_{RA}-1}\alpha^{(p,RA)}_k e^{j2\pi(k+Kk_1+\bar{k})\Delta f_{RA}(t-N_{CP,l}^{RA}T_c - t^{RA}_{start})}
\end{equation}
其中$\alpha^{(p,RA)}_k$是PRACH频域的信号，我们要把它映射到$l$ timeslot上面\\
$t^{\text{RA}}_{\text{start}} \leq t < t^{\text{RA}}_{\text{start}} + (N_{\text{u}} + N_{\text{CP}, l}^{\text{RA}})T_\text{c}$
\subsection{频率部分}
首先我们要确定每个子载波的绝对频率$e^{j2\pi(k+Kk_1+\bar{k})}\Delta f_{RA}$
	\begin{itemize}
	\item $\Delta f_{RA}$是PRACH子载波间隔\\
	1.25kHz或5kHz用于低频段（FR1）的长前导码（$L_{RA}$ = 839），覆盖范围大\\
	15kHz, 30kHz, 60 kHz, 120 kHz：用于高频段（FR2）或特殊场景的短前导码（$L_{RA}$ = 139, 571, 1151），更适合高速和短时延需求。
	\item $\Delta f$是bandwidth part的子载波间隔（inital access bandwidth part 或 the active bandwith part）
	\item $K=\frac{\Delta f}{\Delta f_{RA}}$是比例因子（Scaling Factor）。
	\item $k_1$是一个频率偏移
	\begin{align*}
	k_1 &= k^{\mu}_0 + (N^{\text{start}}_{\text{BWP}, i} - N^{\text{start}, \mu}_{\text{grid}})N^{\text{RB}}_{\text{SC}} - N^{\text{size}, \mu}_{\text{grid}}N^{\text{RB}}_{\text{SC}} + n^{\text{start}}_{\text{RA}}N^{\text{RB}}_{\text{sc}} \\
	&+\begin{cases}
	n_{\text{RA}}N^{\text{RA}}_{\text{RB}}N^{\text{RB}}_{\text{sc}} & L_{\text{RA}}\in\{139, 839\} \\
	n_{\text{RA}}N^{\text{RA}}_{\text{RB}}N^{\text{RB}}_{\text{sc}} & L_{\text{RA}}\in\{571, 1151\} \text{ in FR2-2} \\
	(N^{\text{start}, \mu}_{\text{RB, UL}, n_0+n_{\text{RA}}} - N^{\text{start}, \mu}_{\text{RB, UL}, n_0} ) N^{\text{RB}}_{\text{sc}} & L_{\text{RA}}\in\{571, 1151\} \text{ in FR1}
	\end{cases}
	\end{align*}
		\begin{itemize}
		\item $k^{\mu}_0$是让不同子载波间隔配置（numerology）下偏移（相对最大子载波间隔$\mu_0$对应的资源块的频域中心）
		\[
		k^{\mu}_0 = (N^{\text{start},\mu}_{\text{grid},x} + N^{\text{size},\mu}_{\text{grid},x}/2)N_{\text{sc}}^{\text{RB}} - (N^{\text{start},\mu_0}_{\text{grid},x} + N^{\text{size},\mu_0}_{\text{grid},x}/2)N_{\text{sc}}^{\text{RB}}2^{\mu_0 - \mu}
		\]
		$N^{\text{start},\mu}_{\text{grid},x}$是当前子载波间隔$\mu$下的resouce grid的起始位置（$x$表示DL,UL,SL）\\
		$N^{\text{size},\mu}_{\text{grid},x}$是当前子载波间隔$\mu$下的resouce grid的占用大小（$x$表示DL,UL,SL）\\
		$(N^{\text{start},\mu}_{\text{grid},x} + N^{\text{size},\mu}_{\text{grid},x}/2)N_{\text{sc}}^{\text{RB}}$是当前子载波间隔$\mu$的中心子载波位置\\
		$(N^{\text{start},\mu_0}_{\text{grid},x} + N^{\text{size},\mu_0}_{\text{grid},x}/2)N_{\text{sc}}^{\text{RB}}2^{\mu_0 - \mu}$是最大子载波间隔$\mu_0$的中心子载波位置（换算成$\mu$对应的子载波中心）
		\item $N^{\text{start}}_{\text{BWP}, i}$ 是BWP的最小resource block的编号（initial access或active）\\
		$(N^{\text{start}}_{\text{BWP}, i} - N^{\text{start}, \mu}_{\text{grid}})N^{\text{RB}}_{\text{SC}}$当前上行 BWP 的起始位置，相对于整个资源网格起始位置的偏移量（以子载波为单位）\\
		$(N^{\text{start}}_{\text{BWP}, i} - N^{\text{start}, \mu}_{\text{grid}})N^{\text{RB}}_{\text{SC}} - N^{\text{size}, \mu}_{\text{grid}}N^{\text{RB}}_{\text{SC}}$ 减去半个网格宽度，得到BWP起点相对于网格中心的偏移
		\item $n^{\text{start}}_{\text{RA}}$是相对active BWP的resource block 0的偏移
		\item {\color{darkred}{后面加的部分}}
			\begin{itemize}
			\item $n_{\text{RA}}N^{\text{RA}}_{\text{RB}}N^{\text{RB}}_{\text{sc}}$, $n_{\text{RA}}$表示PRACH在频域的每个传输机会的索引号，$N^{\text{RB}}_{\text{sc}}$表示换算到对应\\
			PUSCH的$\mu$的RB数量。\\
			{\color{darkred}{注意：}}{\color{darkgreen}{FR2-2 PRACH has 139, 571, 1151; FR2-1 PRACH only has 139.}}
			\item 在 FR1 的非授权频谱（例如 5 GHz 频段）中，为了与 Wi-Fi 等其他系统共存，频谱被划分为多个 RB 集合（RB set），每个 RB 集合之间插入了保护带（Guard Band），用来减少干扰。\\
			$N^{\text{start}, \mu}_{\text{RB, UL}, n} $是RB set $n$ 的CRB index\\
			$N^{\text{start}, \mu}_{\text{RB, UL}, n_0} $是RB set $n$ 的第一个PRACH频域机会（$n^{\text{start}}_{RA}$配置）\\
			$N^{\text{start}, \mu}_{\text{RB, UL}, n_0 + n_{RA}} $是RB set $n$ 的第$n_{RA}$个PRACH频域机会\\
			{\color{darkred}{注意：所有$n_{RA}$个PRACH传输机会都在当前RB set内}}
			
			\end{itemize}
		\end{itemize}
	\item $\bar{k}$是PRACH的频率偏移（小）
	\end{itemize}
	
\subsection{时域部分}
$t^{\text{RA}}_{\text{start}} \leq t < t^{\text{RA}}_{\text{start}} + (N_{\text{u}} + N_{\text{CP}, l}^{\text{RA}})T_\text{c}$
\begin{itemize}
\item $N_{\text{u}}$是PRACH 前导码中 Zadoff-Chu 序列（或其它低 PAPR 序列）在时域上所占用的样本数量，{\color{darkred}{具体见TS38211 Table 6.3.3.1-1 \& Table 6.3.3.1-2}}\cite{38211}\\
\item $N_{\text{CP}, l}^{\text{RA}}$是第$l$的OFDM符号上的CP长度，{\color{darkred}{具体见TS38211 Table 6.3.3.1-1 \& Table 6.3.3.1-2}}\cite{38211}\\
\[
N_{\text{CP}, l}^{\text{RA}} = N_{\text{CP}}^{\text{RA}} + n \cdot 16 \kappa
\]
	\begin{itemize}
	\item 当 $\Delta f_{\text{RA}}\in{1.25, 5}$ kHz（即长前导码 839 格式）时：$n=0$，没有额外补偿。
	\item 当 $\Delta f_{\text{RA}}\in{15,30,60,120,480,960}$ kHz（即长前导码 839 格式）时：先看不补偿的PRACH时间$[t^{\text{RA}}_{\text{start}}, t^{\text{RA}}_{\text{start}} + (N_{\text{u}} + N_{\text{CP}}^{\text{RA}})T_\text{c}$ 在一个subframe内和时刻0、时刻$(\Delta f_{\max}N_f/2000)\cdot T_c=0.5$ms重复的次数$n$，每重复一次就要补偿$16 \kappa$。
	\end{itemize}
\item PRACH 的时域位置对齐到了普通上行符号的网格上，便于 gNB 统一管理时序
\[
t^{\text{RA}}_{\text{start}} = t^{\mu}_{\text{start}, l}
\]
\[
t^{\mu}_{\text{start}, l} = 
\begin{cases}
0 & l = 0 \\
t^{\mu}_{\text{start}, l-1} + (N_{\text{u}}^\mu + N_{\text{CP}, l-1}^{\text{RA}})T_\text{c}
\end{cases}
\]
	\begin{itemize}
	\item $N_{\text{u}}^\mu = 2048\kappa \cdot 2^{-\mu}$
	\item $N_{\text{CP}, l}^\mu$
		
	\[
	N_{\text{CP}, l}^\mu = 
	\begin{cases}
	512\kappa \cdot 2^{-\mu} 			& \text{extended CP} \\
	144\kappa \cdot 2^{-\mu} + 16\kappa 	& \text{normal CP, }  l\in \{0, 7\cdot 2^\mu\} \\
	144\kappa \cdot 2^{-\mu} 			& \text{normal CP, }  l\notin \{0, 7\cdot 2^\mu\}  \\
	\end{cases}
	\]
	\item $\Delta f_{\text{RA}} \in\{1.25, 5\}$kHz, $\mu=0$；其他情况$\mu$的值满足$\Delta f_{\text{RA}}=15\cdot 2^\mu$kHz
	\item $l$是prach占用的symbol index
	\[
	l = l_0 + n^{\text{RA}}_tN^{\text{RA}}_{\text{dur}} + 14n^{\text{RA}}_{\text{slot}}
	\]
		\begin{itemize}
		\item $l_0$ 是TS38211 Tables 6.3.3.2-2 to 6.3.3.2-4里面的{\color{darkred}{starting symbol}}\cite{38211}
		\item $n^{\text{RA}}_t=0,\cdots,N_{\text{t}}^{\text{RA, slot}}-1$是PRACH在一个slot内的传输机会。$\Delta L_{\text{RA}} \in\{139, 571, 1151\}$，$N_{\text{t}}^{\text{RA, slot}}$见TS38211 Tables 6.3.3.2-2 to 6.3.3.2-4里面的{\color{darkred}{number of time-domain PRACH occasions within a PRACH slot}}\cite{38211}；$\Delta L_{\text{RA}} = 839$，$N_{\text{t}}^{\text{RA, slot}}=1$\\
		
		{\color{darkred}{TS38211 Tables 6.3.3.2-2 to 6.3.3.2-4 preamble format是A1/B1, A2/B2 or A3/B3时，$n^{\text{RA}}_t$会决定具体的格式}}\cite{38211}
			\begin{itemize}
			\item $n^{\text{RA}}_t=N_{\text{t}}^{\text{RA, slot}}-1$，发B1, B2 or B3
			\item 其他时刻发A1, A2 or A3
			\end{itemize}					
		
		\item $N^{\text{RA}}_{\text{dur}}$见TS38211 Tables 6.3.3.2-2 to 6.3.3.2-4里面的{\color{darkred}{PRACH duration}}\cite{38211}
		
		\item $n^{\text{RA}}_{\text{slot}}$的设置
			\begin{itemize}
			\item $\Delta f_{\text{RA}} \in\{1.25, 5, 15, 60\}$kHz, $n^{\text{RA}}_{\text{slot}}=0$
			\item $\Delta f_{\text{RA}} \in\{30, 120\}$kHz, TS38211 Tables 6.3.3.2-2 to 6.3.3.2-4里面的{\color{darkred}{Number of PRACH slots within a subframe}}或{\color{darkred}{Number of PRACH slots within a 60 kHz slot}}是1，则$n^{\text{RA}}_{\text{slot}}=1$；否则$n^{\text{RA}}_{\text{slot}}\in\{0,1\}$\cite{38211}
			\item $\Delta f_{\text{RA}} \in\{480, 960\}$kHz\\
			TS38211 Tables 6.3.3.2-4里面的{\color{darkred}{Number of PRACH slots within a 60 kHz slot}}是1\\
			$n^{\text{RA}}_{\text{slot}}=7$ for $\Delta f_{\text{RA}}=480$kHz and $n^{\text{RA}}_{\text{slot}}=15$ for $\Delta f_{\text{RA}}=960$kHz
			
			Tables 6.3.3.2-4里面的{\color{darkred}{Number of PRACH slots within a 60 kHz slot}}是2\\
			$n^{\text{RA}}_{\text{slot}}\in\{3,7\}$ for $\Delta f_{\text{RA}}=480$kHz and $n^{\text{RA}}_{\text{slot}}\in\{7,15\}$ for $\Delta f_{\text{RA}}=960$kHz
			\end{itemize}
		
		
		
		\end{itemize}
	\end{itemize}


\end{itemize}

\newpage
\bibliographystyle{IEEEtran}
\bibliography{../../common_latex/import_reference}
\end{document}
